\documentclass[11pt]{article}

\usepackage[a4paper,margin=1in]{geometry}
\usepackage[T2A]{fontenc}
\usepackage[utf8]{inputenc}
\usepackage[russian]{babel}

\usepackage{amsmath,amssymb,amsthm}
\usepackage{mathtools}
\usepackage{hyperref}

\title{Игрушечная модель процентного риска БК в непрерывном времени: FV, дефолт и оптимальное управление}
\author{}
\date{}

\newcommand{\E}{\mathbb{E}}
\newcommand{\Q}{\mathbb{Q}}
\newcommand{\Pp}{\mathbb{P}}
\newcommand{\1}{\mathbf{1}}

\begin{document}
\maketitle

\section{Обзор модели}
Рассматривается упрощённый банковский бизнес: банк непрерывно выдаёт кредиты с фиксированной ставкой (купон фиксируется в момент выдачи) и фондирует их по короткой ставке $r_t$. Процентный риск возникает из-за того, что стоимость фондирования мгновенно реагирует на изменения $r_t$, а средняя доходность кредитного портфеля меняется медленно, поскольку определяется историей выдач.

Работа ведётся в непрерывном времени. Все процессы адаптированы к фильтрации $(\mathcal{F}_t)_{t\ge 0}$ на вероятностном пространстве $(\Omega,\mathcal{F},\Pp)$ (реальная мера) и эквивалентной мере $\Q$ (риск-нейтральная мера). Если не оговорено иное, динамики выписаны под общей мерой $M\in\{\Pp,\Q\}$.

\section{Динамика короткой ставки (Vasicek)}
Короткая ставка $(r_t)_{t\ge 0}$ следует модели Васичека:
\begin{equation}
dr_t = \kappa(\theta^{M}-r_t)\,dt + \sigma\,dW_t^{M},
\qquad r_0 \ \text{задано},
\label{eq:vasicek}
\end{equation}
где $\kappa>0$ --- скорость возврата к среднему, $\sigma>0$ --- волатильность, $\theta^{M}$ --- долгосрочный уровень под мерой $M$.

Далее предполагается, что $\Pp$ и $\Q$ отличаются только параметром долгосрочного уровня:
\[
\theta^{\Pp} \neq \theta^{\Q}, \qquad \kappa,\sigma \ \text{одинаковы под обеими мерами.}
\]

\section{Состояние баланса и управление}
Вводим следующие переменные состояния:
\begin{itemize}
  \item $L_t \ge 0$ --- объём (ноционал) кредитного портфеля (позиция);
  \item $\bar c_t$ --- средняя (book) ставка/купон по outstanding кредитам;
  \item $K_t$ --- капитал, выделенный на бизнес (поглощает убытки от процентного риска);
  \item $u_t \ge 0$ --- управление (интенсивность выдачи).
\end{itemize}

Для простоты амортизацию/выбытие кредитов моделируем \emph{экспоненциальным runoff} с интенсивностью
\[
\lambda := \frac{1}{T},
\]
где $T$ интерпретируется как эффективная средняя жизнь портфеля (например, $T=5$ лет).

\subsection{Динамика объёма позиции}
Пусть абсолютный поток новых выдач (origination) равен
\begin{equation}
q_t := u_t \,\lambda L_t,
\label{eq:qdef}
\end{equation}
то есть $u_t$ --- \emph{относительное} управление:
\begin{itemize}
  \item $u_t=1$: выдачи компенсируют runoff, $L_t$ примерно постоянно;
  \item $u_t<1$: портфель сокращается;
  \item $u_t>1$: портфель растёт.
\end{itemize}
Тогда динамика объёма:
\begin{equation}
dL_t = (q_t - \lambda L_t)\,dt = \lambda L_t (u_t-1)\,dt.
\label{eq:Ldyn}
\end{equation}

\subsection{Ставка новых выдач и динамика средней ставки}
Новые кредиты выдаются по купону
\begin{equation}
c_t^{\mathrm{new}} = r_t + s(u_t),
\label{eq:newcoupon}
\end{equation}
где $s(u)$ --- спред, зависящий от темпа выдачи (например, при агрессивном росте спред может снижаться).

Ключевой момент: динамика средней ставки должна зависеть от \emph{доли нового притока} относительно outstanding. Введём ``купонный номинал''
\[
A_t := L_t\,\bar c_t.
\]
За $dt$:
\begin{itemize}
  \item выбытие $\lambda L_t dt$ уносит средний купон $\bar c_t$ $\Rightarrow$ вклад $-\lambda L_t \bar c_t\,dt$ в $dA_t$;
  \item новые выдачи $q_t dt$ приходят с купоном $c_t^{\mathrm{new}}$ $\Rightarrow$ вклад $+q_t c_t^{\mathrm{new}}\,dt$ в $dA_t$.
\end{itemize}
Отсюда
\[
dA_t = (q_t c_t^{\mathrm{new}} - \lambda L_t \bar c_t)\,dt,
\qquad
dL_t = (q_t - \lambda L_t)\,dt.
\]
Дифференцируя $\bar c_t=A_t/L_t$, получаем
\begin{equation}
d\bar c_t = \frac{q_t}{L_t}\big(c_t^{\mathrm{new}}-\bar c_t\big)\,dt.
\label{eq:cbardyn_general}
\end{equation}
С учётом \eqref{eq:qdef} и \eqref{eq:newcoupon}:
\begin{equation}
d\bar c_t = u_t\,\lambda\big(r_t + s(u_t) - \bar c_t\big)\,dt.
\label{eq:cbardyn}
\end{equation}

Замечание: при $u_t=0$ новых выдач нет, и \eqref{eq:cbardyn} даёт $d\bar c_t=0$ --- средняя ставка портфеля замораживается, что соответствует интуиции.

\subsection{Процентный доход и динамика капитала}
Мгновенная ставка чистого процентного дохода (до кредитного риска и расходов):
\begin{equation}
\pi_t = L_t(\bar c_t - r_t).
\label{eq:profitrate}
\end{equation}
Капитал накапливает прибыль (дивиденды не рассматриваются):
\begin{equation}
dK_t = \pi_t\,dt = L_t(\bar c_t - r_t)\,dt.
\label{eq:Kdyn}
\end{equation}

\section{Справедливая стоимость (FV) на бесконечном горизонте}
Определим справедливую стоимость (economic value) как дисконтированное ожидание будущих потоков под мерой $\Q$:
\begin{equation}
\mathrm{FV}(x)
=
\E^{\Q}\left[\int_{0}^{\infty}
\exp\!\left(-\int_{0}^{t} r_s\,ds\right)\,
\pi_t\,dt \ \Big|\ x_0=x\right],
\label{eq:FV_inf}
\end{equation}
где $x=(r,\bar c,L,K)$ --- начальное состояние.

\section{FV с учётом дефолта и post-default runoff}
Дефолт определяем как исчерпание капитала:
\begin{equation}
\tau := \inf\{t\ge 0:\ K_t \le 0\}.
\label{eq:tau}
\end{equation}

После дефолта бизнес не может выдавать новые кредиты, поэтому вводим правило:
\begin{equation}
u_t = 0 \quad \text{для всех } t \ge \tau.
\label{eq:postdefault_u}
\end{equation}
Тогда портфель после дефолта только сокращается:
\[
dL_t = -\lambda L_t\,dt,\qquad t\ge\tau,
\]
а средняя ставка фиксируется:
\[
d\bar c_t = 0,\qquad t\ge\tau,
\]
поскольку новых выдач нет (см.\ \eqref{eq:cbardyn}).

Справедливая стоимость с дефолтом (и с последующим runoff позиции) задаётся тем же интегралом, но с учётом ограничения \eqref{eq:postdefault_u}:
\begin{equation}
\mathrm{FV}^{\mathrm{def}}(x)
=
\E^{\Q}\left[\int_{0}^{\infty}
\exp\!\left(-\int_{0}^{t} r_s\,ds\right)\,
\pi_t\,dt \ \Big|\ x_0=x,\ u_t=0\ \forall t\ge\tau\right].
\label{eq:FVdef_inf}
\end{equation}
В отличие от ``полной остановки бизнеса'', здесь после дефолта остаются потоки от существующей позиции, пока она закрывается.

\section{Оптимальное управление: постановка, дефолт на горизонте $H$}
Далее рассматриваем управляемую постановку под реальной мерой $\Pp$ с явным ограничением/штрафом на дефолт на конечном горизонте $H>0$. Управление $u_t$ выбирается адаптивно по состоянию.

\subsection{Класс допустимых управлений}
Пусть
\[
\mathcal{U} := \left\{(u_t)_{t\ge 0}:\ u_t\ \text{прогрессивно измеримо, } 0\le u_t \le u_{\max}\right\},
\]
и дополнительно после дефолта действует ограничение $u_t=0$ для $t\ge\tau$.

\subsection{Стационарная value function}
Рассмотрим стационарную (не зависящую явно от времени) value function с субъективной ставкой дисконтирования $\rho\ge 0$ .
\begin{equation}
V(x)
=
\sup_{u\in\mathcal{U}}
\E^{\Pp}\left[
\int_{0}^{\tau}
e^{-\rho t}\,\pi_t\,dt
\;-\;
\eta\,\1_{\{\tau \le H\}}
\ \Big|\ x_0=x
\right].
\label{eq:value_stationary}
\end{equation}
Здесь $\eta\ge 0$ --- штраф за дефолт к горизонту $H$.

\subsection{Пример управления: mean-reversion к целевому уровню капитальной обеспеченности}
Введём показатель ``капитальная обеспеченность'' (обратное плечо)
\begin{equation}
\ell_t := \frac{K_t}{L_t},\qquad (L_t>0),
\label{eq:ell}
\end{equation}
и целевое значение $\ell^\star>0$.

Интуитивная политика: при $\ell_t<\ell^\star$ сокращать выдачи (уменьшать риск), при $\ell_t>\ell^\star$ --- можно наращивать.
Одна из гладких параметрических форм:
\begin{equation}
u_t
=
u_{\min} + (u_{\max}-u_{\min})\,
\sigma\!\left(\alpha(\ell_t-\ell^\star)\right),
\qquad
\sigma(x)=(1+e^{-x})^{-1},
\label{eq:policy_sigmoid}
\end{equation}
где $\alpha>0$ задаёт агрессивность обратной связи.
Эта политика является частным случаем марковского управления $u_t = u(r_t,\bar c_t,L_t,K_t)$.

\section{HJB уравнение (стационарная постановка, до дефолта)}
Ниже выписывается HJB для стационарной задачи на бесконечном горизонте (для простоты записи), где дефолт учтён как граничное условие по капиталу. Это стандартная форма, удобная для анализа структуры оптимального управления.

Рассмотрим стационарную задачу:
\begin{equation}
V(x)
=
\sup_{u\in\mathcal{U}}
\E^{\Pp}\left[
\int_{0}^{\tau}
e^{-\rho t}\,\pi_t\,dt
\ \Big|\ x_0=x
\right],
\qquad
V(r,\bar c,L,0)=0.
\label{eq:stationary_infinite}
\end{equation}
Тогда, при достаточной гладкости $V$, для области $K>0$ выполняется уравнение Беллмана:
\begin{equation}
\rho V
=
\sup_{u\in[0,u_{\max}]}
\Big[
\pi(r,\bar c,L)
+
\mathcal{L}^{u}V
\Big],
\label{eq:HJB}
\end{equation}
где $\pi(r,\bar c,L)=L(\bar c-r)$, а генератор $\mathcal{L}^{u}$ имеет вид
\begin{align}
\mathcal{L}^{u}V
&=
\kappa(\theta^{\Pp}-r)\,V_r
+\frac{1}{2}\sigma^2\,V_{rr}
+\lambda L(u-1)\,V_L
+u\lambda\big(r+s(u)-\bar c\big)\,V_{\bar c}
+L(\bar c-r)\,V_K.
\label{eq:generator}
\end{align}
Граничное условие дефолта:
\begin{equation}
V(r,\bar c,L,K)=0\quad \text{при } K\le 0.
\label{eq:bc_default}
\end{equation}

Замечание: если требуется именно постановка с конечным горизонтом $H$ и штрафом $\eta\,\1_{\{\tau\le H\}}$, то возникает явная зависимость от времени (терминальные условия при $t=H$). Однако для практического моделирования и численных методов часто используют стационарную аппроксимацию \eqref{eq:HJB} (например, с $\rho>0$), а риск дефолта учитывают отдельными ограничениями/штрафами на горизонте $H$ при оценке политики.

\section{Итого: система уравнений}
До дефолта ($K_t>0$) управляемая динамика:
\begin{align}
dr_t &= \kappa(\theta^{\Pp}-r_t)\,dt + \sigma\,dW_t^{\Pp},\\
dL_t &= \lambda L_t(u_t-1)\,dt,\\
d\bar c_t &= u_t\lambda\big(r_t + s(u_t) - \bar c_t\big)\,dt,\\
dK_t &= L_t(\bar c_t - r_t)\,dt,
\end{align}
где $\tau=\inf\{t:K_t\le 0\}$.

После дефолта ($t\ge\tau$) действует ограничение $u_t=0$, поэтому
\[
dL_t=-\lambda L_t\,dt,\qquad d\bar c_t=0,\qquad \pi_t=L_t(\bar c_t-r_t).
\]

\end{document}
